% source: David Mumford, "Abelian Varieties", TIFR Studies in Mathematics 5, Oxford UP, 1970
% pdf_page: 0062
% doc_page: 51 (Chapter II, Section 5, "Cohomology and base change")
% retrieved: 2026-07-11
% transcribed_by: claude-fable-5 (vision, rendered at 200dpi; scanned book, no text layer)

% requires: amsmath (\xrightarrow), amssymb (\mathfrak), mathrsfs (\mathscr)

\DeclareMathOperator{\coker}{coker}

% [running head: ALGEBRAIC THEORY VIA VARIETIES, page 51]

% [conclusion of the statement of Corollary 2, begun on the previous page:
%  conditions (i) and (ii) and the supplementary assertion]

\textit{(i) $y \longmapsto \dim_{k(y)} H^p(X_y, \mathscr{F}_y)$ is a constant
function,}

\textit{(ii) $R^p f_*(\mathscr{F})$ is a locally free sheaf $\mathscr{E}$ on
$Y$, and for all $y \in Y$, the natural map}
\[
  \mathscr{E} \otimes_{\mathcal{O}_Y} k(y) \longrightarrow H^p(X_y, \mathscr{F}_y)
\]
\textit{is an isomorphism.}

\textit{If these conditions are fulfilled, we have further that}
\[
  R^{p-1} f_*(\mathscr{F}) \otimes_{\mathcal{O}_Y} k(y) \longrightarrow
  H^{p-1}(X_y, \mathscr{F}_y)
\]
\textit{is an isomorphism for all $y \in Y$.}

\paragraph{Proof.}
Again assume $Y$ affine, $K^{\cdot}$ as in the proposition.
(ii) $\Rightarrow$ (i) is obvious. To prove (i) $\Rightarrow$ (ii), we need
two lemmas.

\paragraph{Lemma 1.}
\textit{If $Y$ is reduced and $\mathscr{F}$ a coherent sheaf on $Y$ such that
$\dim_{k(y)}[\mathscr{F} \otimes_{\mathcal{O}_Y} k(y)] = r$, all $y \in Y$,
then $\mathscr{F}$ is locally free of rank $r$ on $Y$.}

\paragraph{Proof.}
For any $y \in Y$, let $\sigma_1, \ldots, \sigma_r \in \mathscr{F}_y$ lift
generators of $\mathscr{F}_y \otimes k(y)$. Since
$\sigma_1, \ldots, \sigma_r$ are extendable to sections in a neighborhood of
$y$, we have a homomorphism
$\sigma : \mathcal{O}_Y^r|_V \to \mathscr{F}|_V$ defined in a neighborhood
$V$ of $y$. Then $\sigma$ is surjective on the stalks at $y$, by Nakayama's
lemma, so $\coker(\sigma)$ is zero at $y$ and hence in a neighborhood of $y$.
Thus, we may assume $\sigma$ to be surjective. Then by assumption, for every
$y' \in V$, the map
\[
  \sigma \otimes k(y') : k(y')^r \to \mathscr{F}_{y'} \otimes_{\mathcal{O}_{y'}} k(y')
\]
is an isomorphism. Thus, if $\mathfrak{D}$ is the kernel of $\sigma$, we have
$\mathfrak{D}_{y'} \subset \mathfrak{M}_{y'}\, \mathcal{O}_{y'}^r$ for each
$y' \in V$. Since $Y$ is reduced, this means that $\mathfrak{D} = (0)$. Thus
$\sigma$ is an isomorphism.

We apply this in the following

\paragraph{Lemma 2.}
\textit{Let $Y$ be a reduced, noetherian affine scheme, and let}
\[
  \mathscr{F} \xrightarrow{\ \phi\ } \mathfrak{D}
\]
% [statement of Lemma 2 continues on the next page]
